\practicechapter{第 2 章实践：CPU 如何执行指令}

本章对应第一册第 2 章“CPU 如何执行指令”。正文讲 fetch、decode、execute、寄存器、内存访问、分支和乱序执行。本章不试图模拟 CPU，而是做一个可操作的小目标：从 ELF 的 \code{.text} section 找到机器码区域，再用反汇编把字节、指令和函数符号连起来。

\practicesection{一、对应原书问题：从 \code{.text} 到指令}

第一册正文会反复提醒：CPU 执行的是机器指令，不是 C++ 源码行。实践中最容易断开的地方是中间证据。你看到 \code{add_i32} 的源码，却不知道它在二进制里对应哪个符号；你看到 objdump 输出，却不知道这些指令来自文件的哪个区域。

本章要把三件事连起来：\code{cpu1ee\_cli --sections-csv} 输出 \code{.text} 的 offset 和 size；\cmd{objdump -d} 输出反汇编；symbol table 里的函数名指向某个 section index 和 value。三者合起来，才形成“这个函数的机器指令在当前二进制中如何呈现”的证据链。

\practicesection{二、从特例推到规律：地址不是一种}

特例一：section CSV 里的 \code{offset} 是文件偏移，表示从文件开头数多少字节能到达该 section 内容。特例二：反汇编里函数标签旁边的地址通常是虚拟地址或链接地址，不是文件偏移。特例三：运行时进程启用 PIE 和 ASLR 后，实际装载地址还可能再变化。

由此推出规律：看到一个十六进制数，先问它属于哪种地址空间。文件 offset、ELF virtual address、链接地址、运行时虚拟地址不能混用。本卷解析器第一版只报告 section 的 file offset、section address、symbol value，不负责把运行时地址反推回文件位置。

\practicesection{三、实践任务：定位一段机器码}

先输出 section 表：

\begin{lstlisting}[language=bash]
cpu1ee=books/cpu-volume-1-practice/build/executable-evidence-debug/labs/executable_evidence/cpu1ee_cli
$cpu1ee ./hello-o0 --sections-csv | tee results/hello-o0.sections.csv
objdump -d ./hello-o0 | tee reports/hello-o0.disasm
\end{lstlisting}

在 CSV 中找到 \code{.text} 行，记录 \code{address}、\code{offset}、\code{size}。再在反汇编中找到 \code{add_i32} 或 \code{main}。如果找不到 \code{add_i32}，回到第 5 章实践的优化解释：当前构建可能 inline 了它，也可能符号不可见。为了观察指令，可以临时给函数加 \code{__attribute__((noinline))}，但报告必须写明这是为了保留观察点。

阅读源码中的 \code{SectionSummary} 和 \code{SymbolSummary}，回答：\code{section.address}、\code{section.offset}、\code{symbol.value} 分别表达什么？为什么 \code{symbol.value} 不能直接传给 \code{std::span} 当文件偏移？

\practicesection{四、验收：指令级分析不能跳步}

本章交付物是 \filepath{reports/ch02-instruction-evidence.md}。通过标准如下：

\begin{enumerate}
  \item 报告列出 \code{.text} 的 address、offset、size。
  \item 报告贴出一小段反汇编，并说明它来自哪个函数或标签。
  \item 能解释至少一条指令的操作数是寄存器还是内存访问。
  \item 能说明 file offset、section address、symbol value 的区别。
  \item 不把 CPU 微架构行为当成已经证明的事实；本章只证明文件和反汇编层面的指令证据。
\end{enumerate}

完成本章后，读者应该能从“源码函数”走到“当前二进制里可观察到的机器指令”，而不是凭直觉想象 CPU 在执行什么。
